\subsection{\texorpdfstring{同态 \(\alpha\) 与 \(\beta\)；伴随}{同态 alpha 与 beta；伴随}}

设 A 与 B 为拓扑空间，\(u\colon\mathrm{A}\to\mathrm{B}\) 为连续映射。设 \(\mathcal{G}\) 为 B 上的一个预层。由预层直像的定义，层 \(u_{*}u^{*}\mathcal{G}\) 在 B 的开集 U 上的截面是层 \(u^{*}\mathcal{G}\) 在 A 的开集 \(u^{-1}(\mathrm{U})\) 上的截面，也就是一个 B-态射 \(u^{-1}(\mathrm{U})\to\mathrm{E}_{\mathcal{G}}\)。于是，下述对应定义了一个层的态射 \(\widetilde{\mathcal{G}}\to u_{*}u^{*}\mathcal{G}\)：把截面 \(s\) （\(\mathrm{E}_{\mathcal{G}}\) 在 B 的一个开集 U 上的）对应为截面 \(s\circ u\) （\(\mathrm{E}_{\mathcal{G}}\) 在 \(u^{-1}(\mathrm{U})\) 上的）。该态射与典范态射 \(\sigma_{\mathcal{G}}\colon\mathcal{G}\to\widetilde{\mathcal{G}}\) （I, p. 53）的复合是一个预层态射 \(\mathcal{G}\to u_{*}u^{*}\mathcal{G}\)，将记作 \(\beta_{\mathcal{G}}^{u}\)，甚至记作 \(\beta_{\mathcal{G}}\) （当映射 \(u\) 无歧义时）。

注. --- 1) 设 A、B、C 为拓扑空间，\(u\colon A \to B\)、\(v\colon B \to C\) 为连续映射；令 \(w = v \circ u\)。设 \(\mathcal{G}\) 为 C 上的一个预层。

设 U 为 C 的开集，s 为 \(E_{G}\) 在 U 上的一个截面。于是 \(\beta_{\mathcal{G}}^{v}(s)\) 是截面 \(s \circ v\) （\(E_{G}\) 在 \(v^{-1}(U)\) 上的），而 \(v_{*}(\beta_{v^{*}\mathcal{G}}^{u})(\beta_{\mathcal{G}}^{v}(s))\) 是截面 \(s \circ v \circ u = s \circ w\) （\(E_{G}\) 在 \(u^{-1}(v^{-1}(U)) = w^{-1}(U)\) 上的）。

由此得 \(\beta_{\mathcal{G}}^{w} = v_{*}(\beta_{v^{*}\mathcal{G}}^{u})\circ \beta_{\mathcal{G}}^{v}\)。

\begin{enumerate}
\def\labelenumi{\arabic{enumi})}
\setcounter{enumi}{1}
\tightlist
\item
  若 \(\gamma\colon\mathcal{G}_{1}\to\mathcal{G}_{2}\) 是 B 上的预层态射，则预层态射 \(\beta_{\mathcal{G}_{2}}\circ\gamma\) 与 \(u_{*}u^{*}(\gamma)\circ\beta_{\mathcal{G}_{1}}\) 相等。事实上，若 V 是 B 的开集且 \(s\in\mathcal{G}_{1}(V)\)，则 \(\beta_{\mathcal{G}_{1}}(s)\) 是截面 \(t\) （A \(\times_{B}E_{\mathcal{G}_{1}}\) 在 \(u^{-1}(V)\) 上的），由 \(x\mapsto(x,[V,s,u(x)])\) 定义。于是 \(t\) 在 \(u^{*}(\gamma)\) 下的像是 A \(\times_{B}E_{\mathcal{G}_{2}}\) 在 \(u^{-1}(V)\) 上由 \(x\mapsto(x,[V,\gamma(s),u(x)])\) 给出的截面。由此得 \(u_{*}u^{*}(\gamma)\circ\beta_{\mathcal{G}_{1}}(s)=\beta_{\mathcal{G}_{2}}(\gamma(s))\)。
\end{enumerate}

命题 4. --- 设 A 与 B 为拓扑空间，\(u\colon A\to B\) 为连续映射，\(\mathcal{G}\) 为 B 上的一个预层，\(\mathcal{F}\) 为 A 上的一个层。

对每个预层态射 \(\varphi\colon\mathcal{G}\to u_{*}\mathcal{F}\)，存在唯一的层的态射 \(\psi\colon u^{*}(\mathcal{G})\to\mathcal{F}\) 满足 \(\varphi=u_{*}(\psi)\circ\beta_{\mathcal{G}}\)。

换言之，典范映射

\[
\operatorname{Mor} \left(u ^ {*} (\mathcal {G}), \mathcal {F}\right)\rightarrow \operatorname{Mor} \left(\mathcal {G}, u _ {*} (\mathcal {F})\right), \quad \psi \mapsto u _ {*} (\psi) \circ \beta_ {\mathcal {G}}
\]

是双射。

采用命题 4 的记号，我们有时记 \(\psi = \varphi^{\sharp}\) 与 \(\varphi = \psi^{\flat}\)。

我们来证明命题 4。把注 2 应用于态射 \(\sigma_{\mathcal{G}}: \mathcal{G} \to \widetilde{\mathcal{G}}\)，便知态射 \(\beta_{\mathcal{G}}\) 等于复合

\[
u _ {*} (u ^ {*} (\sigma_ {\mathcal {G}}) ^ {- 1}) \circ \beta_ {\widetilde {\mathcal {G}}} \circ \sigma_ {\mathcal {G}},
\]

其中 \(u^{*}(\sigma_{\mathcal{G}}) \colon u^{*}(\mathcal{G}) \to u^{*}(\widetilde{\mathcal{G}})\) 是 I, p. 58 的注中的典范同构。设 \(\widetilde{\varphi} \colon \widetilde{\mathcal{G}} \to u_{*}\mathcal{F}\) 为唯一的层的态射，满足 \(\widetilde{\varphi} \circ \sigma_{\mathcal{G}} = \varphi\) (I, p. 54, prop. 3)。于是只需证明：存在唯一的层的态射 \(\widetilde{\psi} \colon u^{*}(\mathcal{G}) \to \mathcal{F}\) 满足 \(u_{*}(\widetilde{\psi}) \circ \beta_{\widetilde{\mathcal{G}}} = \widetilde{\varphi}\)。

因此我们可以假设 \(\mathcal{G}\) 是层。为使层的态射 \(\psi: u^{*}(\mathcal{G}) \to \mathcal{F}\) 满足命题 4 的结论，必须且只需：对 B 的每个开集 V 与每个截面 \(t\) （\(\mathrm{E}_{\mathcal{G}}\) 在 V 上的），有

\[
\varphi_{\mathrm{V}}(t) = \psi_{u^{-1}(\mathrm{V})}(t \circ u \mid u^{-1}(\mathrm{V})).\tag{6}
\]

设 \(U_{0}\) 为 A 的一个开子集，\(s_{0}\) 为 \(u^{*}(\mathcal{G})(\mathrm{U}_{0})\) 的一个元素，换言之，是从 \(U_{0}\) 到 \(E_{\mathcal{G}}\) 中的一个 B-态射。设 \(\mathrm{S}(\mathrm{U}_{0}, s_{0})\) 为三元组 \((\mathrm{U}, \mathrm{V}, t)\) 所成的集，其中 U 是 A 的含于 \(U_{0}\) 的开子集，V 是 B 的满足 \(u(\mathrm{U}) \subset \mathrm{V}\) 的开子集，而 t 是 \(E_{\mathcal{G}}\) 在 V 上的一个截面，且满足

\[
t \circ u \mid \mathrm{U} = s _ {0} \mid \mathrm{U}.\tag{7}
\]

若 \(U_{1}\) 与 \(U_{2}\) 是 A 的开子集且 \(U_{1} \subset U_{2}\)，则用 \(f_{U_{1}U_{2}}\) 表示限制映射 \(\mathcal{F}(U_{2}) \to \mathcal{F}(U_{1})\)。于是对每个 \((\mathrm{U}, \mathrm{V}, t) \in \mathrm{S}(\mathrm{U}_{0}, s_{0})\)，我们有关系式

\[
\begin{array}{l} f_{\mathrm{U}\mathrm{U}_{0}}(\psi_{\mathrm{U}_{0}}(s_{0})) = \psi_{\mathrm{U}}(s_{0} \mid \mathrm{U}) \\ \qquad = \psi_{\mathrm{U}}(t \circ u \mid \mathrm{U}) \\ \qquad = f_{\mathrm{U}u^{-1}(\mathrm{V})}(\psi_{u^{-1}(\mathrm{V})}(t \circ u \mid u^{-1}(\mathrm{V}))). \end{array}
\]

因此，若 \(\psi \colon u^{*}(\mathcal{G})\to \mathcal{F}\) 满足 (6)，则我们有

\[
f_{\mathrm{U}\mathrm{U}_{0}}(\psi_{\mathrm{U}_{0}}(s_{0})) = f_{\mathrm{U}u^{-1}(\mathrm{V})}(\varphi_{\mathrm{V}}(t)).\tag{8}
\]

我们来证明：对每一点 \(a\) （属于 \(\mathrm{U}_0\) 者），存在三元组 \((\mathrm{U},\mathrm{V},t)\in \mathrm{S}(\mathrm{U}_0,s_0)\) 使得 \(a\in \mathrm{U}\)。事实上，设 \(a\) 为 \(\mathrm{U}_0\) 的一点。存在 B 中包含 \(u(a)\) 的开邻域 V，以及截面 \(t\) （平展空间 \(\mathbf{E}_{\mathcal{G}}\) 在 V 上的），使得 \(t(u(a)) = s_0(a)\) (I, p. 33, prop. 9)。令 \(\mathrm{U}_1 = u^{-1} (\mathrm{V})\cap \mathrm{U}_0\)。截面 \(s_0\mid \mathrm{U}_1\) 与 \(t\circ u\mid \mathrm{U}_1\) ——即 \(\mathrm{U}_1\) 上的平展空间 \(\mathbf{E}_{\mathcal{G}}\times_{\mathrm{B}}\mathrm{U}_1\) 的截面——在点 \(a\) 处相合。由 I, p. 34 的 prop. 11, b)，二者相合的点所成的集是 \(\mathrm{U}_1\) 的一个包含 \(a\) 的开集 U。于是三元组 \((\mathrm{U},\mathrm{V},t)\) 属于 \(\mathrm{S}(\mathrm{U}_0,s_0)\)。

于是，公式 (8) 与层的性质 \((\mathrm{F}_{1})\) (I, p. 43) 蕴涵 \(\psi\) 的唯一性。

设 (U, V, t) 与 \((\mathrm{U}^{\prime}, \mathrm{V}^{\prime}, t^{\prime})\) 为 \(\mathrm{S}(\mathrm{U}_{0}, s_{0})\) 的元素。由关系式 (7)，t 与 \(t^{\prime}\) 在 \(u(\mathrm{U} \cap \mathrm{U}^{\prime})\) 上的限制相合。由 I, p. 34 的 prop. 11, b)，存在 B 的开集 W 满足 \(u(\mathrm{U} \cap \mathrm{U}^{\prime}) \subset \mathrm{W} \subset \mathrm{V} \cap \mathrm{V}^{\prime}\) 且 \(t \mid W = t^{\prime} \mid W\)。因此我们有

\[
f_{u^{-1}(\mathrm{W})u^{-1}(\mathrm{V})}(\varphi_{\mathrm{V}}(t)) = \varphi_{\mathrm{W}}(t\mid \mathrm{W}) = \varphi_{\mathrm{W}}(t^{\prime}\mid \mathrm{W}) = f_{u^{-1}(\mathrm{W})u^{-1}(\mathrm{V}^{\prime})}(\varphi_{\mathrm{V}^{\prime}}(t^{\prime}))
\]

从而

\[
f_{(\mathrm{U}\cap\mathrm{U}^{\prime})u^{-1}(\mathrm{V})}(\varphi_{\mathrm{V}}(t)) = f_{(\mathrm{U}\cap\mathrm{U}^{\prime})u^{-1}(\mathrm{V}^{\prime})}(\varphi_{\mathrm{V}^{\prime}}(t^{\prime})).\tag{9}
\]

由层 F 的性质 \((\mathrm{F}_{1})\) 与 \((\mathrm{F}_{2})\)，存在唯一的元素 \(s'\) （\(\mathcal{F}(\mathrm{U}_{0})\) 的），使得对每个三元组 \((\mathrm{U},\mathrm{V},t)\) （取自

\(\mathrm{S}(\mathrm{U}_{0},s_{0})\) 者），有：

\[
f_{\mathrm{U}\mathrm{U}_{0}}(s^{\prime}) = f_{\mathrm{U}u^{-1}(\mathrm{V})}(\varphi_{\mathrm{V}}(t)).\tag{10}
\]

把这个元素记作 \(\psi_{\mathrm{U}_{0}}(s_{0})\)。

设 \(U_{1}\) 为含于 \(U_{0}\) 的开集，且 \(s_{1}=s_{0}|U_{1}\)。若 \((\mathrm{U},\mathrm{V},t)\in\mathrm{S}(\mathrm{U}_{1},s_{1})\)，则 U 是含于 \(U_{0}\) 的开集且 \(t\circ u|U=s_{1}|U=s_{0}|U\)，因此 \((\mathrm{U},\mathrm{V},t)\in\mathrm{S}(\mathrm{U}_{0},s_{0})\)，而关系式 (10) 这时蕴涵

\[
f_{\mathrm{U}\mathrm{U}_{1}}(f_{\mathrm{U}_{1}\mathrm{U}_{0}}(\psi_{\mathrm{U}_{0}}(s_{0}))) = f_{\mathrm{U}\mathrm{U}_{0}}(\psi_{\mathrm{U}_{0}}(s_{0})) = f_{\mathrm{U}u^{-1}(\mathrm{V})}(\varphi_{\mathrm{V}}(t)).
\]

由 \(\psi_{\mathrm{U}_{1}}(s_{1})\) 的定义，我们于是有 \(\psi_{\mathrm{U}_{1}}(s_{1}) = f_{\mathrm{U}_{1}\mathrm{U}_{0}}(\psi_{\mathrm{U}_{0}}(s_{0}))\)。这就证明了族 \(\psi = (\psi_{\mathrm{U}})\) 是从 \(u^{*}(\mathcal{G})\) 到 \(\mathcal{F}\) 中的层的态射。

我们来证明 \(\psi\) 满足关系式 (6)。设 V 为 B 的开子集，\(t\) 为 \(\mathrm{E}_{\mathcal{G}}\) 在 V 上的一个截面。若 \(\mathrm{U} = u^{-1}(\mathrm{V})\) 且 \(s = t \circ u \mid \mathrm{U}\)，则三元组 \((\mathrm{U}, \mathrm{V}, t)\) 属于 \(\mathrm{S}(\mathrm{U}, s)\)，而关系式 (6) 是关系式 (10) （应用到 \(\mathrm{U} = \mathrm{U}_0\)）的直接推论。

命题 5. — 设 A 与 B 为拓扑空间，u: A → B 为连续映射。

\begin{enumerate}
\def\labelenumi{\alph{enumi})}
\tightlist
\item
  设 \(\mathcal{G}_1\) 与 \(\mathcal{G}_2\) 为 B 上的预层，\(\gamma\colon\mathcal{G}_1\to\mathcal{G}_2\) 为预层的一个态射。再设 \(\mathcal{F}\) 为 A 上的层，\(\varphi\colon\mathcal{G}_2\to u_*\mathcal{F}\) 为预层的一个态射。我们有等式
\end{enumerate}

\[
(\varphi \circ \gamma) ^ {\sharp} = \varphi^ {\sharp} \circ u ^ {*} (\gamma).\tag{11}
\]

\begin{enumerate}
\def\labelenumi{\alph{enumi})}
\setcounter{enumi}{1}
\tightlist
\item
  设 \(\mathcal{F}_1\) 与 \(\mathcal{F}_2\) 为 A 上的层，\(\mathcal{G}\) 为 B 上的一个预层。设 \(\varphi\colon\mathcal{F}_1\to\mathcal{F}_2\) 为层的一个态射，\(\gamma\colon\mathcal{G}\to u_*\mathcal{F}_1\) 为预层的一个态射。我们有关系式
\end{enumerate}

\[
(u _ {*} (\varphi) \circ \gamma) ^ {\sharp} = \varphi \circ \gamma^ {\sharp}.
\]

\begin{enumerate}
\def\labelenumi{\alph{enumi})}
\tightlist
\item
  由 \(\varphi^{\sharp}\) 的定义与 I, p. 60 的注 2，我们有
\end{enumerate}

\[
\varphi \circ \gamma = u _ {*} (\varphi^ {\sharp}) \circ \beta_ {\mathcal {G} _ {2}} \circ \gamma = u _ {*} (\varphi^ {\sharp}) \circ u _ {*} u ^ {*} (\gamma) \circ \beta_ {\mathcal {G} _ {1}}.
\]

因此，\(\varphi \circ \gamma = u_{*}(\varphi^{\sharp}\circ u^{*}(\gamma))\circ \beta_{\mathcal{G}_{1}}\)；由此得关系式 (11)。

\begin{enumerate}
\def\labelenumi{\alph{enumi})}
\setcounter{enumi}{1}
\tightlist
\item
  由 \(\gamma^{\sharp}\) 的定义，我们有
\end{enumerate}

\[
u _ {*} (\varphi) \circ \gamma = u _ {*} (\varphi) \circ u _ {*} (\gamma^ {\sharp}) \circ \beta_ {\mathcal {G}} = u _ {*} (\varphi \circ \gamma^ {\sharp}) \circ \beta_ {\mathcal {G}},
\]

从而得到所宣布的关系式，这里用到了 \((u_{*}(\varphi)\circ\gamma)^{\sharp}\) 的定义。

令 \(\alpha_{\mathcal{F}} = \mathrm{Id}_{u_*(\mathcal{F})}^\sharp\)；它是唯一的层的态射 \(\rho \colon u^*(u_*(\mathcal{F})) \to \mathcal{F}\)，满足

\[
\operatorname{Id} _ {u _ {*} (\mathcal {F})} = u _ {*} (\rho) \circ \beta_ {u _ {*} (\mathcal {F})}.
\]

把关系式 (11) 应用于 \(\mathcal{G}_2 = u_*(\mathcal{F})\) 与态射 \(\varphi = \mathrm{Id}_{u_*(\mathcal{F})}\)，就对每个预层态射 \(\gamma\colon\mathcal{G}\to u_*(\mathcal{F})\) 给出分解式

\[
\gamma^ {\sharp} = \alpha_ {\mathcal {F}} \circ u ^ {*} (\gamma).
\]

于是由命题 4 推出：对每个层的态射 \(\psi\) （从 \(u^{*}(\mathcal{G})\) 到 \(\mathcal{F}\) 中的），\(\psi^{\flat}\) 是唯一的态射 \(\varphi\colon\mathcal{G}\to u_{*}(\mathcal{F})\) 满足 \(\psi=\alpha_{\mathcal{F}}\circ u^{*}(\varphi)\)。

例. --- 1) 考察拓扑空间 B、B 的子空间 A，并用 \(i\colon A \to B\) 表示典范单射。设 \((E,p)\) 与 \((E',p')\) 为 B-空间。取 \(\mathcal{G}\) 为层 \(\underline{\mathcal{M}or}_{B}(E;E')\) （I, p. 45, 例 4），取 \(\mathcal{F}\) 为层 \(\underline{\mathcal{M}or}_{A}(E_{A};E_{A}')\)。对 B 的每个开集 V 与每个 V-态射 \(f\colon E_{V} \to E_{V}'\)，令 \(\varphi_{V}(f) = f_{V \cap A}\)，其中 \(f_{V \cap A}\) 是 \((V \cap A)\) -态射 （从 \(E_{V \cap A}\) 到 \(E_{V \cap A}'\) 中的、由 \(f\) 诱导者）。如此定义的族 \(\varphi = (\varphi_{V})\) 是从 \(\mathcal{G}\) 到 \(i_*\mathcal{F}\) 中的层的态射。由命题 4，存在唯一的层的态射 \(\psi\colon \underline{\mathcal{M}or}_{B}(E;E')_{A} \to \underline{\mathcal{M}or}_{A}(E_{A};E_{A}')\)，使得有

\[
\psi_ {\mathrm{V} \cap \mathrm{A}} (\sigma_ {\mathcal {G}} (\mathrm{V}, f) \mid \mathrm{V} \cap \mathrm{A}) = f _ {\mathrm{V} \cap \mathrm{A}}\tag{12}
\]

对 B 的每个开集 V 与每个 \(f \in \mathcal{C}_{\mathrm{V}}(\mathrm{E}_{\mathrm{V}}; \mathrm{E}_{\mathrm{V}}^{\prime})\) 成立。态射 \(\psi\) 称为从 \(\underline{\mathcal{M}or_{B}}(\mathrm{E}; \mathrm{E}^{\prime})_{\mathrm{A}}\) 到 \(\underline{\mathcal{M}or_{A}}(\mathrm{E}_{\mathrm{A}}; \mathrm{E}_{\mathrm{A}}^{\prime})\) 中的典范态射。

若 A 缩为一个单点 \(a\)，则该态射等同于从 \(a\) 处的茎（层 \(\underline{\mathcal{M}or}_{\mathrm{B}}(\mathrm{E};\mathrm{E}^{\prime})\) 的）到集合 \(\mathcal{C}(\mathrm{E}_a;\mathrm{E}_a')\) 中的态射。

通过过渡到子层，态射 \(\psi\) 诱导出从 \(\mathcal{I}som_{\mathrm{B}}(\mathrm{E};\mathrm{E}')_{\mathrm{A}}\) 到 \(\mathcal{I}som_{\mathrm{A}}(\mathrm{E}_{\mathrm{A}};\mathrm{E}_{\mathrm{A}}^{\prime})\) 中的一个典范态射。

\begin{enumerate}
\def\labelenumi{\arabic{enumi})}
\setcounter{enumi}{1}
\tightlist
\item
  设 A、B、C 为拓扑空间，\(u\colon A\to B\)、\(v\colon B\to C\) 为连续映射；令 \(w=v\circ u\)。设 E 与 E' 为 C-空间。从 \(\underline{\mathcal{M}or}_{C}(E;E')_{A}\) 到 \(\underline{\mathcal{M}or}_{A}(E_{A};E'_{A})\) 中的典范态射是这样一个态射与另一个典范态射的复合：前者为 \(\underline{\mathcal{M}or}_{C}(E;E')_{A}\to\underline{\mathcal{M}or}_{B}(E_{B};E'_{B})_{A}\) （由 \(\underline{\mathcal{M}or}_{C}(E;E')_{B}\) 到 \(\underline{\mathcal{M}or}_{B}(E_{B};E'_{B})\) 中的典范态射导出者），后者为从 \(\underline{\mathcal{M}or}_{B}(E_{B},E'_{B})_{A}\) 到 \(\underline{\mathcal{M}or}_{A}(E_{A},E'_{A})\) 中的典范态射。
\end{enumerate}
% semantic-fragments: legacy-input-seam
\relax{}
